\paragraph{Spectral positivity of the exact correlator.}
Suppose the matrix dynamics admits a self-adjoint Hamiltonian \(H\)
bounded below with a normalizable ground state \(\Omega\), and let
\(T_{\mu\nu}\) be a self-adjoint composite rank-two operator, the
fluctuation of the composite metric to which the external metric
couples through the dictionary. Then
\(\langle\Omega|T_{\mu\nu}(t)\,T_{\rho\sigma}(0)|\Omega\rangle
=\sum_n e^{-\mathrm i(E_n-E_0)t}
\langle\Omega|T_{\mu\nu}|n\rangle\langle n|T_{\rho\sigma}|\Omega
\rangle\) over the spectral resolution of \(H\): the measure is
positive and carried by squared matrix elements, and the residue at
any isolated simple real pole of the spin-two-projected two-point
function is a sum of squared matrix elements and so non-negative.
Under these hypotheses the exact composite spin-two two-point
function has no negative-residue real pole at any energy; the extra
pole of the fixed-\(L\) form and of the leading-logarithmic \(F\)
above is a property of those truncated kernels. These hypotheses
are not established in the present hypothesis class: a self-adjoint
Hamiltonian bounded below with a normalizable ground state is the
content of Assumption~\ref{ass:lorentzian}; the Euclidean measures
of Section~\ref{sec:status-E} fail the positivity of
\cite{OsterwalderSchrader1973,OsterwalderSchrader1975} on the
invariant polynomial algebra at finite \(N\), and the tree-level
reflected form at the Euclidean flat point of
Section~\ref{sec:status-D} is negative on the gauge-invariant
translation probes; neither the fluctuation class nor \(T_{\mu\nu}\)
is derived (Problem~D). Even under the hypotheses the sublemma
would not by itself discharge Assumption~\ref{ass:EFT-a2}: it
constrains the exact kernel, whereas the pole statements of this
section concern computed truncations whose remainder is not bounded
near the would-be pole. A two-oscillator quantum-mechanical model
with two positive-norm modes exhibits the artifact: the exact
two-point function of the observed mode has positive residues at
both normal modes, while the quadratic truncation of the same
kernel acquires one real negative-residue pole, located above the
scale of the truncation.

\paragraph{The named completion class and the remainder bound.}
The leading-logarithmic \(F\) is the second-order curvature kernel
of a Barvinsky--Vilkovisky-class heat-kernel expansion
\cite{BarvinskyVilkovisky1990}: about a flat background the
second-order kernel is the only part of that expansion that
contributes to the two-point function, and a completion that fixes
the computed logarithm enters additively,
\(F(x)\mapsto F(x)+R(x)\), with \(R(x)=(x/\Upsilon)r(x)\) and
\(r(0)=0\) in the massless case. Every analytic \(O(x)\) completion
reparametrizes \(\delta_W\), since
\(F(x;\delta_W)+\varepsilon x=F(x;\delta_W-\Upsilon\varepsilon)\);
this is the sense in which finite local terms are absorbed. No
stability statement for the zero-free disc can rest on the
logarithmic coefficient alone: at \(\Upsilon=10\) and
\(\delta_W=0\), the completion \(-c\,x^{2}\) carries no logarithm,
so the \(O(p^4)\) logarithmic coefficient and \(\delta_W\) are
unchanged, and for \(c>1\) it places a zero at \(0<x<1\)
(\(x\approx0.909\) at \(c=1.2\)). What is available is
conditional: with \(m(\Upsilon,\delta_W)\) the minimum of \(|F|\)
over the boundary of the slit disc
\(\{|x|<1\}\setminus(-\infty,0]\) (at \(\Upsilon=10\) and
\(\delta_W=0\), \(m=0.82\), attained at \(|x|=1\)), if
\(\sup|R|\) over the same boundary is below
\(m(\Upsilon,\delta_W)\), then \(F+R\) and \(F\) have the same
number of zeros in the slit disc, so within the parameter region
recorded above \(F+R\) has no zero with \(|x|<1\). The bound on
\(R\) requires the spectral data of the computed form factors on
the unit circle and is not established; it is an input of
Problem~F.
